% What one `tofu apply` reads and writes: configuration, state, real
% infrastructure through the provider, the plan in between.
\documentclass[border=6pt]{standalone}
\usepackage{tofu}
\begin{document}
\begin{tikzpicture}[papier,
    box/.style={bloc=#1, minimum width=36mm, minimum height=13mm, text width=33mm, align=center}]
  \node[box=Pvert] (conf) {Configuration\\[-1pt]{\scriptsize\ttfamily *.tf}};
  \node[box=Pbleu, below=14mm of conf] (state) {State\\[-1pt]{\scriptsize\ttfamily terraform.tfstate}};
  \node[box=Plavande, right=22mm of conf, yshift=-13.5mm, minimum height=22mm] (core)
    {\ttfamily tofu\\[2pt]{\sffamily\scriptsize construit le graphe,\\ compare, planifie}};
  \node[box=Porange, right=20mm of core, yshift=13.5mm] (plan) {Plan\\[-1pt]{\scriptsize\ttfamily + \ \textasciitilde{} \ -/+ \ -}};
  \node[box=Psable, right=20mm of core, yshift=-13.5mm] (prov) {Provider\\[-1pt]{\scriptsize plugin : docker, aws...}};
  \node[box=Pjaune, right=18mm of prov] (real) {Infrastructure\\ réelle};
  \draw[fl] (conf.east) -- node[et, above, sloped] {lit} (core.north west);
  \draw[fl] (state.east) -- node[et, below, sloped] {lit} (core.south west);
  \draw[fl] (core.north east) -- node[et, above, sloped] {produit} (plan.west);
  \draw[flbi] (core.south east) -- node[et, below, sloped] {gRPC} (prov.west);
  \draw[flbi] (prov) -- node[et, above] {API} node[et, below] {lit, crée} (real);
  \draw[fld=Pbrun] (core.south) -- ++(0,-12mm) -| node[etb=Pbrun, pos=0.25, below] {après apply, tofu écrit le nouvel état} (state.south);
\end{tikzpicture}
\end{document}
